\subsubsection{設計根拠}

\term{設計根拠}{Design Rationale} は、なぜその設計判断をしたかを説明する情報である。設計根拠には、前提、検討した代替案、採用基準、トレードオフ、却下した案の理由が含まれる。

設計チームにとって、現在の判断理由は明らかに見えることが多い。しかし、数か月後、数年後、または担当者交代後には、その理由が失われる。設計根拠を記録すれば、保守担当者は「なぜこの構造なのか」「なぜ別案を採らなかったのか」を理解しやすい。

却下した判断を記録することも有用である。前提や制約が変わったとき、以前は却下された案が有効になる場合がある。逆に、過去に明確な理由で却下された案を、理由を忘れて再び採用してしまうことを防げる。

\begin{pointbox}{実務の見方}
設計根拠は、巨大な説明文である必要はない。「判断」「理由」「代替案」「影響」「未解決事項」を短く残すだけでも、後の保守で大きな価値を持つ。
\end{pointbox}
